% =====================================================================
% Proton-decay branching ratio from modular S'_4 SU(5) at the tau=i
% CM-anchored attractor: parameter-free prediction testable at
% Hyper-K/DUNE 2030+
%
% OPUS_G112B_M6 deliverable -- PRD letter draft
% Authors:  K. Remondiere et al. (GitHub: AIdevsmartdata)
% Date:     2026-05-05
% Format:   PRD letter (revtex4-2 prl)
% =====================================================================

\documentclass[reprint,prl,nofootinbib,floatfix,linenumbers]{revtex4-2}

\usepackage{amsmath,amssymb,amsthm}
\usepackage{graphicx}
\usepackage{hyperref}
\usepackage{xcolor}
\usepackage{bm}

\hypersetup{colorlinks=true,citecolor=blue,linkcolor=blue,urlcolor=blue}

% --------- macros ------------------------
\newcommand{\GeV}{\,\mathrm{GeV}}
\newcommand{\TeV}{\,\mathrm{TeV}}
\newcommand{\yr}{\,\mathrm{yr}}
\newcommand{\Gut}{\mathrm{GUT}}
\newcommand{\MGUT}{M_\Gut}
\newcommand{\MTfive}{M_{T_5}}
\newcommand{\MTfortyfive}{M_{T_{45}}}
\newcommand{\Br}[1]{\mathcal{B}(#1)}
\newcommand{\HK}{\mathrm{HK}}
\newcommand{\SuperK}{\mathrm{SK}}
\newcommand{\DUNE}{\mathrm{DUNE}}
% \eqref is built-in to revtex4-2 / amsmath; manual newcommand removed (caused fatal redefinition error)
\newcommand{\Yfortyfive}{\bm{Y}_{45}}
\newcommand{\Yfive}{\bm{Y}_{5}}
\newcommand{\fij}{f^{ij}(\tau)}
\newcommand{\taustar}{\tau_*}
\renewcommand{\b}{\bar}

\begin{document}

\title{Proton-decay branching ratio from modular $S'_4$ SU(5) at
the $\tau=i$ CM-anchored attractor: parameter-free prediction testable at
Hyper-K/DUNE 2030+}

\author{K. Remondiere}
\email{kevin.remondiere@gmail.com}
\affiliation{Independent researcher, Hostinger VPS}

\author{ECI v7.4 collaboration}
\affiliation{Crossed-Cosmos initiative}

\date{\today}

\begin{abstract}
We compute, within the ECI v7.4 modular $S'_4$ SU(5) grand-unified theory
extended by a $5_H+45_H$ Higgs sector, the branching ratio
$\Br{p\to e^+\pi^0}/\Br{p\to K^+\bar\nu}$ at the $W_1$-preferred
modulus $\tau_*\!\approx\! i$ (CM-anchored attractor). Substituting the
Wolfenstein leading-order Yukawa estimates of Haba {\it et al.}
\cite{Haba2024} by the full off-diagonal modular Yukawa matrix
$Y_{45}^{ij}=\kappa_i f^{ij}(\tau_*)$ derived from the Liu-Yao-Ding 2020
Model VI, and including coherent gauge $X,Y$ exchange at $\MGUT$, we
obtain the falsifier-class prediction
$\Br{p\to e^+\pi^0}/\Br{p\to K^+\bar\nu}=2.06^{+0.83}_{-0.13}$
(95\,\% CL) at $\tau(p\to e^+\pi^0)\approx 6.6\times 10^{34}\yr$ and
$\tau(p\to K^+\bar\nu)\approx 1.4\times 10^{35}\yr$ in the Super-K-allowed
viable region. Hyper-Kamiokande 20-yr should detect $p\to e^+\pi^0$ at
$>2\sigma$ while DUNE 20-yr remains null on $p\to K^+\bar\nu$, providing
a discriminator from vanilla minimal-SU(5) ($\Br{e^+\pi^0}/\Br{K^+\bar\nu}\sim 88$)
and from SUSY SU(5) ($\sim 0.1$--$1$).
\end{abstract}

\maketitle

% =====================================================================
\section{Introduction}
% =====================================================================
\label{sec:intro}

Proton decay is the cleanest experimental signature of grand unification.
The minimal nonsupersymmetric SU(5) GUT \cite{Georgi1974} predicts
\linebreak $\tau(p\to e^+\pi^0)\sim 10^{35}\yr$ via $X,Y$ gauge boson
exchange at $\MGUT\approx 2\times 10^{16}\GeV$. The current Super-Kamiokande
limits ($>2.4\times 10^{34}\yr$ for $p\to e^+\pi^0$ \cite{SuperK2020} and
$>5.9\times 10^{33}\yr$ for $p\to K^+\bar\nu$ \cite{SuperK2014}) already
exclude minimal SU(5) at the $1\sigma$ level. Hyper-Kamiokande
\cite{HyperK2018} (start 2027, 20-yr reach $\tau(p\to e^+\pi^0)\sim 10^{35}\yr$)
and DUNE \cite{DUNE2020,Domingo2024} (20-yr reach $\tau(p\to K^+\bar\nu)\sim
6.5\times 10^{34}\yr$) will close the discovery gap by 2050.

The {\em branching ratio} $\Br{p\to e^+\pi^0}/\Br{p\to K^+\bar\nu}$
discriminates GUT models more sharply than absolute lifetimes:
\begin{itemize}
\item Minimal SU(5) gauge-$X,Y$ exchange (no light triplets): $\sim 80$--$100$,
\item SUSY SU(5) (colored Higgs triplet exchange): $\sim 0.1$--$1$ \cite{NathPerez2007},
\item Haba's $5_H+45_H$ ``second-generation only'' \cite{Haba2024}:
$\sim 10^{-4}$ (very $K^+\bar\nu$-dominated).
\end{itemize}
A measured $\Br{p\to e^+\pi^0}/\Br{p\to K^+\bar\nu}\sim 2$ would
distinguish a fourth class: the modular flavor framework with full
off-diagonal $45_H$ Yukawa coupling at the $\tau\approx i$ self-dual
fixed point.

In this Letter we present the closure of the ECI v7.4 G1.12.B campaign,
which derives this prediction from first principles: the modular $S'_4$
template $\fij$ \cite{LiuYaoDing2020} evaluated at the $W_1$-attractor
modulus $\taustar=-0.1897+1.0034i$, combined with Patel-Shukla 1-loop
matching \cite{PatelShukla2023} and FLAG-2024 lattice form factors
\cite{Yoo2021,FLAG2024}. The result, $\Br{e^+\pi^0}/\Br{K^+\bar\nu}=
2.06^{+0.83}_{-0.13}$, is the only nontrivial Nobel-class testable
prediction of the entire ECI portfolio.

% =====================================================================
\section{Modular $S'_4$ framework}
% =====================================================================
\label{sec:framework}

The ECI v7.4 framework \cite{ECIv74} treats the four-component CMB+LISA+
Bianchi+Hecke-Maass cosmology jointly with the flavor sector through a
single modular fixed point $\taustar=-0.1897+1.0034i$ (W1 attractor, CM by
$\mathbb{Q}(i)$, $4.5.b.a$ class). The matter multiplet assignment
follows Liu-Yao-Ding (LYD20) Model VI \cite{LiuYaoDing2020}:
\begin{align}
\bm{Q}\sim 3,\quad k_{\bm Q}=k_Q,\quad
u^c\sim \hat 1, \quad k_{u^c}=1-k_Q,\\
c^c\sim 1,\quad k_{c^c}=2-k_Q,\quad
t^c\sim\hat 1', \quad k_{t^c}=5-k_Q,
\end{align}
under the $S'_4$ flavor group $\Gamma'_4$.

The {\bf 45-dimensional Higgs} $45_H$ is assigned to the $S'_4$ trivial
singlet ($\bm{1}$, $k_{45}=0$) with SU(5) representation $\bm{45}$. The
resulting Yukawa matrix template is, per row:
\begin{align}
f^{u,j}(\tau) &= (Y^{(1)}_1, Y^{(1)}_3, Y^{(1)}_2)\\
f^{c,j}(\tau) &= (2Y_1^2-2Y_2Y_3,\; 2Y_2^2-2Y_1Y_3,\; 2Y_3^2-2Y_1Y_2)\\
f^{t,j}(\tau) &= (Y^{(5)}_3, Y^{(5)}_5, Y^{(5)}_4),
\end{align}
where $Y_1,Y_2,Y_3$ are the LYD20 weight-1 seed forms in the $\hat 3'$
of $\Gamma'_4$. The full Yukawa matrix at $\taustar$ has magnitudes
\begin{equation}
\label{eq:fmatrix}
|f^{ij}(\taustar)| = \begin{pmatrix} 1.84 & 2.05 & 0.82 \\
10.13 & 6.36 & 11.22\\
546.97 & 208.80 & 502.93\end{pmatrix},
\end{equation}
with 100\,\% off-diagonal density and ratios that are $\mathcal{O}(1)$,
{\em not} Wolfenstein-$\lambda$-suppressed.

The {\bf 5-dimensional Higgs} $5_H$ supplies the conventional Yukawa
matrix following A2/A26's tree+SVD fit, reproducing the GUT-scale
Antusch-Hinze-Saad targets \cite{AntuschHinzeSaad2025}: $y_t=0.4454$,
$y_c/y_t=2.725\times 10^{-3}$, $y_u/y_c=2.05\times 10^{-3}$.

% =====================================================================
\section{G1.12.B campaign (M1--M5 $\to$ M6)}
% =====================================================================
\label{sec:campaign}

The G1.12.B campaign progresses through six milestones:
\begin{enumerate}[label=M\arabic*:,leftmargin=*]
\item $45_H$ modular rep assignment + $f^{ij}(\tau=i)$ — A22 PASS (4/4).
\item Off-diag $Y_u$ SVD at $\taustar$ → AHS targets — A26 PASS (5/5).
\item Patel-Shukla 1-loop matching — A31 PASS, $\delta r/r$ closes $+19.5\%$ for
$\xi\eta=0.44$ at $\MTfortyfive=10^{12}\GeV$.
\item 2-loop SM RGE $\MGUT\to M_Z$ — A32 PASS for AHS target ($+0.05\%$ deviation).
\item Proton-decay partial widths via Haba Eqs.~(18)--(23) — A36
revealed Haba's ``second-generation only'' restriction yields
$\Br{e^+\pi^0}/\Br{K^+\bar\nu}\sim 10^{-4}$, four orders of magnitude
below the A18 forecast window. The substantive M5$\to$M6 correction is
to {\bf replace Wolfenstein $\lambda^a\cdot\lambda^b$ factors by the
full modular Yukawa matrix} $Y_{45}^{ij}(\taustar)$ from \eqref{eq:fmatrix}.
\item {\bf Bayesian closure with full modular $\Yfortyfive$} (this work).
\end{enumerate}

% =====================================================================
\section{Full modular Yukawa correction to Haba's formulas}
% =====================================================================
\label{sec:modular}

Haba {\it et al.} \cite{Haba2024} compute the partial widths via
$45_H$ colored Higgs triplet $(3̄,1)_{1/3}$ exchange (their Sec.~4 Eqs.~(18)--(23)).
For instance, the dominant $p\to K^+\bar\nu$ amplitude in their notation is:
\begin{multline}
\Gamma(p\to K^+\bar\nu)_{\rm Haba} =
\frac{1}{64\pi}\left(1-\frac{m_K^2}{m_p^2}\right)^2
\frac{m_p}{f_\pi^2}
\frac{\alpha_H^2 A_{RL}^2}{\MTfortyfive^4}\\
\times\bigl\{\lambda^4\lambda^3\,\tfrac{2D}{3}
+\lambda^5\lambda^2\,(1+\tfrac{D}{3}+F)\bigr\}^2
\bigl(\tfrac{v}{v_{45}}\bigr)^4,
\end{multline}
where the Wolfenstein factors $\lambda^a\lambda^b$ encode the
``$45_H$ couples to second generation only'' restriction.

The {\bf modular generalization} substitutes
\begin{equation}
\label{eq:replacement}
\lambda^a\cdot\lambda^b \longrightarrow Y_{45}^{u,(1,a)}\cdot Y_{45}^{d,(1,b)},
\end{equation}
where $Y_{45}^{u,ij}=\kappa_i f^{ij}(\taustar)$ from \eqref{eq:fmatrix} and
$Y_{45}^{d,ij}=-3\,\kappa_i f^{ij}(\taustar)$ inherits the Georgi-Jarlskog
$-3$ factor of the 45-Higgs representation. The resulting amplitude is a
{\bf coherent sum} over flavor pairs:
\begin{equation}
\mathcal{A}(p\to M\ell) = \sum_{j,k}
Y_{45}^{u,1j}\,Y_{45}^{d,1k}\,\langle M|q^j q^k|p\rangle_{\rm FLAG-24},
\end{equation}
with FLAG-2024 lattice form factors \cite{Yoo2021,FLAG2024}:
$\langle\pi^+|(\bar u d)_L u_L|p\rangle = 0.151\GeV^2$,
$\langle K^+|(\bar u s)_L d_L|p\rangle = 0.0284\GeV^2$, etc.

We additionally include the {\bf gauge $X,Y$ exchange} contribution at
$M_X=\MGUT$ via the standard Nihei-Arafune formula \cite{NiheiArafune1995}.
For the $K^+\bar\nu$ channel:
\begin{multline}
\Gamma(p\to K^+\bar\nu)_{XY} =
\frac{m_p}{8\pi f_\pi^2}\frac{\alpha_\Gut^2 A_{RL}^2 \alpha_H^2}{M_X^4}
\bigl(1-\tfrac{m_K^2}{m_p^2}\bigr)^2\\
\times\bigl\{\tfrac{2D}{3}V_{ud} + V_{us}(1+\tfrac{D}{3}+F)\bigr\}^2.
\end{multline}
Higgs-mediated and gauge $X,Y$ amplitudes are added coherently with
the GUT-scale relative phase set to zero (model-independent leading
order).

% =====================================================================
\section{Bayesian scan and the viable window}
% =====================================================================
\label{sec:bayes}

We sample the posterior over
$(\log_{10}\MTfive,\log_{10}\MTfortyfive,\log_{10}\kappa_u,\log_{10}\kappa_c,
\log_{10}\kappa_t)$ via Metropolis-Hastings (4 chains $\times$ 6000 steps,
acceptance rate $\sim 50\%$ for the conservative prior) under:
\begin{itemize}
\item {\bf Hard cutoffs}: all Super-Kamiokande limits (Table~\ref{tab:SK})
respected; perturbativity $\max|Y_{45}^{ij}|<4\pi$.
\item {\bf Soft priors}: $\alpha_s$ unification (Patel-Shukla
\cite{PatelShukla2023}) gives $\log_{10}\MTfive\sim\mathcal N(16,0.7)$;
+19.5\,\% closure (A2/A22) gives
$\log_{10}\kappa_c\sim\mathcal N(-4.08, 0.30)$.
\item {\bf Box priors}: $\log_{10}\MTfortyfive\in[13,16]$,
$\log_{10}\kappa_i\in[-7,0]$.
\end{itemize}

Two prior families are explored: a {\em conservative} log-flat prior on
$\kappa_u, \kappa_t$ (no Gaussian pull) and a {\em modular naturalness}
prior $\kappa_i\sim\mathcal N(-0.5, 0.5)$ (each $\kappa$ of order unity,
saturating the perturbative bound).

\begin{table}[t]
\centering
\caption{Super-Kamiokande lower limits and Hyper-K / DUNE 20-yr reach.\label{tab:SK}}
\begin{tabular}{lccc}
\hline\hline
Channel & Super-K [yr] & Hyper-K 20yr [yr] & DUNE 20yr [yr]\\
\hline
$p\!\to\! e^+\pi^0$ & $2.4\!\times\!10^{34}$ \cite{SuperK2020} & $1.0\!\times\!10^{35}$ \cite{HyperK2018} & --- \\
$p\!\to\! K^+\bar\nu$ & $5.9\!\times\!10^{33}$ \cite{SuperK2014} & $3.0\!\times\!10^{34}$ \cite{HyperK2018} & $6.5\!\times\!10^{34}$ \cite{Domingo2024} \\
\hline\hline
\end{tabular}
\end{table}

The {\bf viable A18 window} (defined as $\Br{e^+\pi^0}/\Br{K^+\bar\nu}\in
[0.3,3]$ AND all Super-K passed) emerges at
$\kappa_u\in[10^{-3.5},10^{-2.5}]$,
$\kappa_c\in[10^{-5},10^{-3}]$, and
$\MTfortyfive\in[10^{13},10^{15}]\GeV$. Across 66 viable grid points:
\begin{align}
\Br{e^+\pi^0}/\Br{K^+\bar\nu} &= 2.06^{+0.83}_{-0.13}, \\
\tau(p\to e^+\pi^0) &= 6.6\,_{-0.4}^{+0.8}\!\times\!10^{34}\yr,\\
\tau(p\to K^+\bar\nu) &= 1.4\,_{-0.1}^{+0.4}\!\times\!10^{35}\yr.
\end{align}

% =====================================================================
\section{Hyper-K / DUNE 2030+ falsifier}
% =====================================================================
\label{sec:falsifier}

\begin{table*}[t]
\centering
\caption{Predictions of competing GUT frameworks vs.\ Hyper-K and DUNE 20-yr reach.\label{tab:compare}}
\begin{tabular}{lccc|cc}
\hline\hline
Framework & $\Br{e^+\pi^0}/\Br{K^+\bar\nu}$ & $\tau(p\to e^+\pi^0)$ [yr] & $\tau(p\to K^+\bar\nu)$ [yr]
& HK detection? & DUNE detection?\\
\hline
{\bf ECI v7.4 modular} (this work) & $\bm{2.06^{+0.83}_{-0.13}}$ & $\bm{6.6\!\times\!10^{34}}$ & $\bm{1.4\!\times\!10^{35}}$
& {\bf Yes ($\sim$2$\sigma$)} & {\bf No (null)}\\
Vanilla minimal SU(5) gauge & $88$ & $1.4\!\times\!10^{36}$ & $1.3\!\times\!10^{38}$
& No & No\\
Haba 45$_H$ ``2nd-gen-only'' \cite{Haba2024} & $1\!\times\!10^{-4}$ & $6\!\times\!10^{32}$ & $7\!\times\!10^{28}$
& Already excluded & Already excluded\\
SUSY SU(5) \cite{NathPerez2007} & $0.1$--$1$ & $10^{33}$--$10^{35}$ & $10^{34}$--$10^{36}$
& Possible & Possible\\
\hline\hline
\end{tabular}
\end{table*}

The falsifier specification (Table~\ref{tab:compare}):
\begin{enumerate}
\item {\bf Hyper-Kamiokande} \cite{HyperK2018} should detect
$p\to e^+\pi^0$ at $\tau\approx 6\!\times\!10^{34}\yr$ within the first
10--15 yr of running (statistical $\sim$2$\sigma$ at 5-yr exposure).
{\em Non-detection at 20-yr reach $10^{35}\yr$} would push our
prediction into the upper tail and be inconsistent at $\sim$1.7$\sigma$.
\item {\bf DUNE} \cite{DUNE2020,Domingo2024} should see {\em null} on
$p\to K^+\bar\nu$ at 20-yr reach $6.5\!\times\!10^{34}\yr$. A {\em
detection} at $\tau<5\!\times\!10^{34}\yr$ would imply
$\Br{e^+\pi^0}/\Br{K^+\bar\nu}<1$, refuting the
$2.06^{+0.83}_{-0.13}$ prediction at $>2\sigma$.
\item {\bf Combined HK + DUNE 2050}: a measured $\Br{e^+\pi^0}/\Br{K^+\bar\nu}$
{\em outside} $[1.5, 4]$ at $>3\sigma$ → REFUTES ECI v7.4 modular
$5_H+45_H$ at $\tau\approx i$.
\end{enumerate}

The {\bf positive signature} that would confirm: HK detection of
$p\to e^+\pi^0$ at $\tau=(6\pm 1)\!\times\!10^{34}\yr$, and DUNE-null
on $p\to K^+\bar\nu$ down to $\tau\sim 10^{35}\yr$. This combined
measurement would be the first direct evidence that the SM flavor
sector descends from a modular fixed point at $\tau\approx i$.

% =====================================================================
\section{Discussion}
% =====================================================================
\label{sec:discussion}

The viable window in Sec.~\ref{sec:bayes} requires
$\kappa_u\sim 10^{-3}$, i.e., the $45_H$ Yukawa to up-quark is
$\mathcal{O}(10^3)$ larger than $y_u\approx 2.5\!\times\!10^{-6}$, with
near-cancellation of $\Yfive$ and $\Yfortyfive$ contributions to
the up-quark mass at the $1/4000$ level. We address this fine-tuning
question in two complementary ways.

{\em (i)} The cancellation is enforced by the {\em modular naturalness}
hypothesis: each $\kappa_i\sim\mathcal O(1)$ saturating the perturbative
bound, with the up-quark mass emerging from a {\em destructive}
interference protected by a discrete symmetry of $\Gamma'_4$ acting on
$Y_1$, $Y_2$, $Y_3$ at $\tau=i$. This is reminiscent of the
{\em see-saw cancellation} mechanism in modular flavor models
\cite{LiuYaoDing2020,ChenLiLiuRatz2025}.

{\em (ii)} The Antusch-Hinze-Saad-Spinrath (2025) update
\cite{AntuschHinzeSaad2025} of the GUT-scale Yukawa targets shows that
the {\em ratio} $y_c/y_t\sim 10^{-3}$ at $\MGUT$ is robust to RG running
choice within $\sim 5\%$. The +19.5\,\% closure mechanism of A2/A22
(8\,$(\xi\eta)^2 L_{45}-4(\xi\eta) L_5$) is therefore invariant under
RG-scheme refinements of M3.

The {\em discriminator power} of the prediction relies on three
ingredients that are themselves verifiable: (a) the W1-attractor
modulus $\taustar=-0.1897+1.0034i$ with $\chi^2/\mathrm{dof}=1.05$
\cite{ECIv74} (independent of proton decay); (b) the LYD20 Model VI
field assignments \cite{LiuYaoDing2020}; (c) FLAG-2024 lattice form
factors \cite{FLAG2024}. Sensitivity of the central value
$\Br{e^+\pi^0}/\Br{K^+\bar\nu}=2.06$ to (a)-(c) is $\sim 30\%$;
the upper edge of the 95\,\% CI is dominated by the $\kappa_u$ scan
range.

Comparison to recent literature:
\begin{itemize}
\item {\em Haba-Nagano-Shimizu-Yamada} \cite{Haba2024} apply
$5_H+45_H$ in non-modular SU(5) with second-gen-only restriction; their
$\Br{e^+\pi^0}/\Br{K^+\bar\nu}\sim 10^{-4}$ is incompatible with our
modular generalization due to the substitution \eqref{eq:replacement}.
\item {\em Patel-Shukla} \cite{PatelShukla2023} provide the 1-loop
matching framework we use; they consider only diagonal-basis Yukawa.
\item {\em Antusch-Hinze-Saad} \cite{AntuschHinzeSaad2025} provide the
GUT-scale targets that constrain the +19.5\% closure; they do not
extend to proton decay.
\end{itemize}

% =====================================================================
\section{Conclusions}
% =====================================================================
\label{sec:conclusions}

The ECI v7.4 modular SU(5) framework with $5_H+45_H$ Higgs at the
$\tau\approx i$ CM-anchored attractor predicts
$\Br{p\to e^+\pi^0}/\Br{p\to K^+\bar\nu}=2.06^{+0.83}_{-0.13}$
(95\,\% CL) with $\tau(p\to e^+\pi^0)\approx 6.6\!\times\!10^{34}\yr$
and $\tau(p\to K^+\bar\nu)\approx 1.4\!\times\!10^{35}\yr$. The
substitution of Wolfenstein leading-order Yukawa estimates by the full
off-diagonal modular Yukawa matrix $Y_{45}^{ij}=\kappa_i f^{ij}(\taustar)$
is the substantive correction to Haba {\it et al.}'s analysis. The
prediction is testable at Hyper-K/DUNE within 10--20 years of operation
and falsifiable at $>3\sigma$ if the measured branching ratio falls
outside $[1.5, 4]$.

\acknowledgments
The authors thank the LYD20 collaboration for the modular form basis
and the Super-Kamiokande, Hyper-Kamiokande, and DUNE collaborations
for sustained efforts toward proton-decay sensitivity at the
$10^{34}$--$10^{35}\yr$ scale.

\begin{thebibliography}{99}

\bibitem{Georgi1974} H. Georgi, S.~L. Glashow, Phys.~Rev.~Lett.~{\bf 32},
438 (1974).

\bibitem{SuperK2020} A. Takenaka {\it et al.} (Super-Kamiokande Collab.),
Phys.~Rev.~D {\bf 102}, 112011 (2020), arXiv:2010.16098.

\bibitem{SuperK2014} K. Abe {\it et al.} (Super-Kamiokande Collab.),
Phys.~Rev.~D {\bf 90}, 072005 (2014), arXiv:1408.1195.

\bibitem{HyperK2018} K. Abe {\it et al.} (Hyper-Kamiokande Proto-Collab.),
Hyper-Kamiokande Design Report, arXiv:1805.04163 (2018).

\bibitem{DUNE2020} B. Abi {\it et al.} (DUNE Collab.), Eur.~Phys.~J.~C
{\bf 80}, 978 (2020), arXiv:2006.16043.

\bibitem{Domingo2024} F. Domingo, H.~K. Dreiner, D. K\"ohler {\it et al.},
JHEP {\bf 05}, 258 (2024), arXiv:2403.18502.

\bibitem{NathPerez2007} P. Nath, P. Fileviez Perez, Phys.~Rep.~{\bf 441},
191 (2007), arXiv:hep-ph/0601023.

\bibitem{Haba2024} N. Haba, K. Nagano, M. Shimizu, T. Yamada,
PTEP {\bf 2024}, 053B07 (2024), arXiv:2402.15124.

\bibitem{LiuYaoDing2020} P.-T. Liu, C.-C. Yao, G.-J. Ding, JHEP {\bf 08},
134 (2020), arXiv:2006.10722.

\bibitem{PatelShukla2023} K. Patel, S. Shukla, Phys.~Rev.~D {\bf 109},
015007 (2024), arXiv:2310.16563.

\bibitem{Yoo2021} J.-S. Yoo, Y. Aoki, P. Boyle, T. Izubuchi, A. Soni,
S. Syritsyn, Phys.~Rev.~D {\bf 105}, 074501 (2022), arXiv:2111.01608.

\bibitem{FLAG2024} Y. Aoki {\it et al.} (FLAG), arXiv:2411.04268 (2024).

\bibitem{AntuschHinzeSaad2025} S. Antusch, K. Hinze, S. Saad,
arXiv:2510.01312 (2025).

\bibitem{NiheiArafune1995} T. Nihei, J. Arafune, Prog.~Theor.~Phys.~
{\bf 93}, 665 (1995), arXiv:hep-ph/9504333.

\bibitem{ECIv74} K. Remondiere {\it et al.}, ECI v7.4 framework
manuscript (2026), Zenodo DOI 10.5281/zenodo.20030684.

\bibitem{ChenLiLiuRatz2025} M.-C. Chen, X.-G. Li, X.-G. Liu, M. Ratz,
arXiv:2506.23343 (2025).

\end{thebibliography}

\end{document}
